% !TeX root = ../../Book1.tex
% 纲要标题：### 2.3 正规子群与商群
% 纲要位置：计划纲要.md，第 173 行。

\section{正规子群与商群}
\label{b1:sec:0203}

% 纲要内容（仅作写作计划，尚非教材正文）：
%
% * 正规性及商群构造
% * 商群的泛性质
% * 单群概念与简单例子
%

陪集给出了集合的划分。要使代表元的乘法下降为陪集集合上的运算，须检验换代表元后输出是否仍是同一个陪集；在群中，这个相容条件恰好由正规性描述。

\ProseParagraphBreak
本节先建立商群，再说明哪些同态能够通过商映射。单群的定义放在最后，它所排除的是非平凡真正规子群，而不是所有非平凡真子群。

\subsection{正规性及商群构造}
\label{b1:subsec:0203:01}

对任意子集 \(A,B\subseteq G\)，集合乘积 \(AB=\{ab\mid a\in A,\ b\in B\}\) 总有定义；两个陪集的集合乘积未必是单个陪集。这里要检验的不是集合乘积能否形成，而是候选规则
\[
(aH)\star(bH)\coloneqq(ab)H
\]
能否在陪集集合上良定义。为此，换取同一输入陪集的代表元时，输出必须不变。

\ProseParagraphBreak
暂将子群记为 \(N\leq G\)。在 \(aN,bN\) 中分别换取代表 \(a'=an,b'=bm\)，其中 \(n,m\in N\)，乘积就成为 \(anbm\)。要使它仍落在 \(abN\) 中，需要把夹在 \(a,b\) 之间的 \(n\) 移到 \(b\) 右侧。利用
\[
nb=b(b^{-1}nb),\qquad
a'b'=anbm=ab(b^{-1}nb)m,
\]
便可重新写成 \(ab\) 后面乘一个元素。此时需要 \(b^{-1}nb\in N\)，才能保证最后两个因子的积仍在 \(N\) 中。让 \(b\) 遍历 \(G\)、\(n\) 遍历 \(N\)，便得到任意共轭都应保持这个子群的要求。

\begin{definition}\label{b1:def:normal-subgroup}
设 \(G\) 为群，\(N\leq G\)。如果对每个 \(g\in G\) 都有 \(gNg^{-1}=N\)，则将 \(N\) 称为 \(G\) 的\term[zhenggui-ziqun]{正规子群}[normal subgroup]，记为 \(N\trianglelefteq G\)。这里 \(gNg^{-1}\) 表示所有 \(gng^{-1}\ (n\in N)\) 的集合；映射 \(n\mapsto gng^{-1}\) 称为由 \(g\) 给出的\term[gonge]{共轭}[conjugation]。
\end{definition}

\symidx{\(N\trianglelefteq G\)：正规子群关系}

\begin{proposition}\label{b1:prop:normal-criteria}
设 \(G\) 为群，\(N\leq G\)。以下条件等价：\(N\trianglelefteq G\)；对每个 \(g\in G\) 有 \(gNg^{-1}\subseteq N\)；对每个 \(g\in G\) 有 \(gN=Ng\)。
\end{proposition}

\begin{proof}
等号自然给出包含。若所有共轭都满足包含，把 \(g\) 换成 \(g^{-1}\) 得 \(g^{-1}Ng\subseteq N\)。将这个包含式两侧同时左乘 \(g\)、右乘 \(g^{-1}\)，就是把其中每个元素 \(x\) 变为 \(gxg^{-1}\)，包含关系仍然成立。因此
\[
N=g(g^{-1}Ng)g^{-1}\subseteq gNg^{-1}.
\]
逐元素地说，对任意 \(n\in N\)，已知 \(g^{-1}ng\in N\)，所以 \(n=g(g^{-1}ng)g^{-1}\in gNg^{-1}\)。这与原来的 \(gNg^{-1}\subseteq N\) 合起来给出等号。等式 \(gNg^{-1}=N\) 两边右乘 \(g\)，恰为 \(gN=Ng\)；反过来，将 \(gN=Ng\) 两边右乘 \(g^{-1}\) 就恢复共轭等式。
\end{proof}

\begin{proposition}[正规子群的基本例子]\label{b1:prop:normal-basic-examples}
设 \(f:G\rightarrow K\) 为群同态，则 \(\ker f\trianglelefteq G\)。如果群 \(G\) 交换，那么 \(G\) 的每个子群都正规。对任意群 \(G\)，指数为 \(2\) 的每个子群也都正规。
\end{proposition}

\begin{proof}
设 \(f:G\rightarrow K\) 为同态。由\cref{b1:prop:kernel-fibers}，\(\ker f\) 已是子群。若 \(n\in\ker f\)，则
\(f(gng^{-1})=f(g)f(n)f(g)^{-1}=1\)，故 \(g(\ker f)g^{-1}\subseteq\ker f\)；\cref{b1:prop:normal-criteria}给出正规性。若 \(G\) 交换，则对任意子群 \(H\)、\(h\in H\) 和 \(g\in G\)，有 \(ghg^{-1}=h\)，所以 \(H\) 正规。

最后设 \([G:H]=2\)。由\cref{b1:prop:coset-equivalence}，左、右陪集各有两个。若 \(g\in H\)，则 \(gH=Hg=H\)；若 \(g\notin H\)，则两者都是各自划分中除 \(H\) 外的唯一一块，因而都等于 \(G\setminus H\)。所以对所有 \(g\) 有 \(gH=Hg\)，再用\cref{b1:prop:normal-criteria}即得正规性。
\end{proof}

\begin{theorem}[商群构造]\label{b1:thm:quotient-group}
设 \(G\) 为群，\(N\leq G\)。候选规则 \((aN)\star(bN)=(ab)N\) 在左陪集集合上良定义，当且仅当 \(N\trianglelefteq G\)。满足此条件时，该运算使 \(G/N\) 成为群，将所得群称为 \(G\) 关于 \(N\) 的\term[shangqun]{商群}[quotient group]；其单位元为 \(N\)，逆元公式为 \((aN)^{-1}=a^{-1}N\)。商映射 \(q:G\rightarrow G/N\)、\(q(g)=gN\) 是满同态，核恰为 \(N\)。
\end{theorem}

\begin{proof}
若 \(N\) 正规，写 \(a'=an\)、\(b'=bm\)，其中 \(n,m\in N\)，则
\[
a'b'=anbm=ab(b^{-1}nb)m.
\]
正规性保证最后两个因子之积属于 \(N\)，故 \(a'b'N=abN\)，证明良定义。结合律由 \((ab)c=a(bc)\) 下降，单位和逆元由
\(1a=a1=a\)、\(aa^{-1}=a^{-1}a=1\) 下降。商映射保持乘法且满射；\(q(g)=N\) 当且仅当 \(g\in N\)。

反过来，假设候选运算 \(\star\) 良定义。它必须对所有代表元更换都给出相同结果，因此为了推出必要条件，可以特别只更换第一项的代表：对任意 \(n\in N\)、\(g\in G\)，由于 \(nN=1N\)，有
\[
(ng)N=(nN)\star(gN)=(1N)\star(gN)=gN.
\]
\cref{b1:prop:coset-equivalence}的陪集相等判据给出 \(g^{-1}ng\in N\)。这个特别选择已经迫使任意 \(n,g\) 满足所需的共轭条件，即对每个 \(g\in G\) 都有 \(g^{-1}Ng\subseteq N\)；再用\cref{b1:prop:normal-criteria}即得 \(N\) 正规。
\end{proof}

正规时，陪集的集合乘积确实等于上述单个陪集：
\[
(aN)(bN)=\{anbm\mid n,m\in N\}=(ab)N.
\]
此后省略临时记号 \(\star\)，以并置表示商群乘法；不正规时，并置仅表示始终有定义的集合乘积，不应误作商群乘法。

\ProseParagraphBreak
商群中的单位元是整个子群 \(N\)，并非单个元素 \(1\)。例如
\(\mathbb Z/4\mathbb Z\) 的 \([2]+[2]=[0]\) 表示两个陪集相加得到单位陪集，而不是说整数 \(2+2=0\)。保持原群等式与商群等式的区别，是使用商结构时最基本的习惯。

\ProseParagraphBreak
非正规子群会怎样破坏乘法，可以沿用三点置换群中的例子检验。在 \(S_3\) 中，令 \(a\) 交换 \(1,2\)，令 \(b\) 交换 \(2,3\)，取 \(H=\{1,a\}\)。由于 \(1H=aH\)，同一个陪集的代表可以取 \(1\)，也可以取 \(a\)。与 \(bH\) 相乘时，候选规则分别给出
\[
(1H)\star(bH)=(1b)H=bH,\qquad
(aH)\star(bH)=(ab)H.
\]
然而 \(b^{-1}ab=bab\) 交换 \(1,3\) 而固定 \(2\)，不属于 \(H\)，故\cref{b1:prop:coset-equivalence}说明 \(abH\neq bH\)。输入陪集没有改变，输出却变了，因而这条规则不能定义陪集集合上的运算。商群定理中的正规性条件正是为了使这种更换代表元的计算结果保持一致。

\ProseParagraphBreak
同一群中也有可以取商的子群。令 \(r=ab\)，直接计算得 \(r\) 是三轮换、\(r^3=1\)，且 \(r^2=ba\)。于是 \(N=\langle r\rangle=\{1,r,r^2\}\) 是三元素子群；\(S_3\) 的六个元素为 \(1,r,r^2,a,b,aba\)，故 \([S_3:N]=2\)。由\cref{b1:prop:normal-basic-examples}，\(N\trianglelefteq S_3\)。商群 \(S_3/N\) 的元素是 \(N\) 与 \(aN=\{a,b,aba\}\)，而 \(a^2=1\) 给出 \((aN)^2=N\)。其乘法表为
\[
\begin{array}{c|cc}
\cdot & N & aN\\ \hline
N & N & aN\\
aN & aN & N
\end{array}
\]
这里 \(H=\{1,a\}\) 不正规，候选运算不能下降；\(N\) 正规，因而确实得到一个二元素商群。

\begin{corollary}[群的同余关系与正规子群]\label{b1:cor:group-congruences}
设 \(G\) 为群。\(G\) 上与乘法相容的等价关系，恰好由正规子群给出。具体地，若 \(\sim\) 是\cref{b1:def:congruence}意义下的同余关系，令 \(N=[1_G]\)，则 \(N\trianglelefteq G\)，且
\[
a\sim b\quad\Longleftrightarrow\quad b^{-1}a\in N,
\qquad [a]_{\sim}=aN
\]
对所有 \(a,b\in G\) 成立。反过来，每个正规子群 \(N\) 都按此公式确定一个同余关系，其单位元所在的类为 \(N\)。
\end{corollary}

\begin{proof}
由自反性，\(1\in N\)。若 \(n,m\in N\)，则 \(n\sim1\)、\(m\sim1\)，乘法相容性给出 \(nm\sim1\)，所以 \(nm\in N\)。在 \(n\sim1\) 两侧左乘 \(n^{-1}\)，得 \(1\sim n^{-1}\)，再用对称性知 \(n^{-1}\in N\)。因此 \(N\) 是子群。对任意 \(g\in G\)，在 \(n\sim1\) 两侧分别左乘 \(g\)、右乘 \(g^{-1}\)，得 \(gng^{-1}\sim1\)；由\cref{b1:prop:normal-criteria}，\(N\) 正规。

左乘 \(b^{-1}\) 把 \(a\sim b\) 变为 \(b^{-1}a\sim1\)，左乘 \(b\) 又能恢复原关系，故得到所述等价式。于是对任意 \(x\in G\)，
\[
x\in[a]_{\sim}\quad\Longleftrightarrow\quad
a^{-1}x\in N\quad\Longleftrightarrow\quad x\in aN,
\]
从而 \([a]_{\sim}=aN\)。反过来，若 \(N\) 正规，\cref{b1:thm:quotient-group}给出商同态 \(q:G\rightarrow G/N\)。关系 \(q(a)=q(b)\) 是等价关系，并由同态性质与乘法相容；\cref{b1:prop:coset-equivalence}说明它正由所给公式描述。令 \(b=1\)，可知单位元所在的类恰为 \(N\)。
\end{proof}

因此，第一章的同余关系方法在群中可以用正规子群完整表达：确定单位元所在的类 \([1]=N\) 后，其余各类就是 \([a]=aN\)，不必再分别指定它们。

\subsection{商群的泛性质}
\label{b1:subsec:0203:02}

设 \(N\trianglelefteq G\)。商映射 \(q:G\rightarrow G/N\) 已经把每个 \(n\in N\) 与单位元视为相同。如果同态 \(f:G\rightarrow K\) 要经过这个商，那么 \(f\) 也必须看不出这些区别，即对每个 \(n\in N\) 都有 \(f(n)=1_K\)。这正是
\[
N\subseteq\ker f.
\]
包含的方向表示：\(q\) 已经忘掉的信息，\(f\) 不能再需要；\(f\) 还可以合并更多元素，所以它的核允许比 \(N\) 更大。下一定理说明，这个条件也足以给出唯一的分解。

\begin{theorem}[商群的泛性质]\label{b1:thm:quotient-universal}
设 \(G,K\) 为群，\(N\trianglelefteq G\)，\(q:G\rightarrow G/N\) 为商映射。对于群同态 \(f:G\rightarrow K\)，存在群同态 \(\overline f:G/N\rightarrow K\) 满足 \(\overline f q=f\)，当且仅当 \(N\subseteq\ker f\)。存在时，\(\overline f\) 唯一，且由 \(\overline f(gN)=f(g)\) 给出。
\end{theorem}

\begin{proof}
如果已有分解，\(n\in N\) 时 \(q(n)=N\) 为商群单位元，故 \(f(n)=1\)。反过来，若 \(N\subseteq\ker f\)，则
由\cref{b1:prop:coset-equivalence}，\(gN=hN\) 蕴含 \(h^{-1}g\in N\)，于是 \(f(h)^{-1}f(g)=1\)，即 \(f(g)=f(h)\)，证明所给公式良定义。由
\(\overline f\bigl((gN)(hN)\bigr)=f(gh)=f(g)f(h)\)，它为同态。每个商群元素都是某个 \(q(g)\)，所以分解等式强制决定其全部取值，给出唯一性。
\end{proof}

当 \(N\subseteq\ker f\) 时，这一分解可以画成交换三角形；虚线表示唯一的群同态。
\[
\begin{tikzcd}[column sep=3em,row sep=2.5em]
G\arrow[r,"q"]\arrow[dr,"f"']&G/N\arrow[d,dashed,"\exists!\,\overline f"]\\
&K
\end{tikzcd}
\]
例如映射 \(\mathbb Z\rightarrow\mathbb Z/3\mathbb Z\)，\(a\mapsto[a]\)，能够经由 \(\mathbb Z/6\mathbb Z\) 分解，因为 \(6\mathbb Z\subseteq3\mathbb Z\)。诱导映射是把模 \(6\) 的类再取模 \(3\)。同一个公式却不能定义从 \(\mathbb Z/4\mathbb Z\) 到 \(\mathbb Z/3\mathbb Z\) 的映射：\(0,4\) 在前者代表同一类，在后者代表不同类。泛性质把这种代表元检验概括为核的包含条件。

\ProseParagraphBreak
无论 \(f\) 是否满射，诱导映射的像与核分别为
\[
\operatorname{im}\overline f=\operatorname{im}f,\qquad
\ker\overline f=\{gN\mid g\in\ker f\}=(\ker f)/N.
\]
事实上，每个商群元素都可写为 \(gN\)，而 \(\overline f(gN)=f(g)\)，所以
\[
\operatorname{im}\overline f
=\{\,\overline f(gN)\mid g\in G\,\}
=\{\,f(g)\mid g\in G\,\}=\operatorname{im}f.
\]
同一取值公式还给出
\[
gN\in\ker\overline f\quad\Longleftrightarrow\quad
f(g)=1_K\quad\Longleftrightarrow\quad g\in\ker f.
\]
因而核恰由 \(\ker f\) 中元素的 \(N\)-陪集组成。这里 \(N\subseteq\ker f\)，且 \(N\) 在 \(G\) 中正规，所以它在子群 \(\ker f\) 中也正规；这些陪集确实组成商群 \((\ker f)/N\)。

\ProseParagraphBreak
若 \(f\) 满射，\(\overline f\) 也满射，但它仍未必单射。因此应区分“商掉一个被核包含的正规子群”与“恰好商掉整个核”。下一节的第一同构定理使用后者，正是为了使诱导映射没有剩余的非平凡核。

\subsection{单群概念与简单例子}
\label{b1:subsec:0203:03}

\begin{definition}\label{b1:def:simple-group}
设 \(G\) 为非平凡群。如果 \(G\) 仅有 \(\{1_G\}\) 和 \(G\) 两个正规子群，则将 \(G\) 称为\term[danqun]{单群}[simple group]。
\end{definition}

对正规子群 \(N\trianglelefteq G\)，两个端点给出 \(G/\{1\}\) 和 \(G/G\)：前者的商映射是同构，后者是只有单位元的群。这里所说的“非平凡真商群”，指由 \(\{1\}\ne N\ne G\) 给出的商 \(G/N\)。此时商映射的核非平凡，故它不单射；同时 \(N\ne G\)，商群又不平凡。因此对非平凡群 \(G\) 有
\[
G\;\text{是单群}\quad\Longleftrightarrow\quad
G\;\text{没有非平凡真商群}.
\]
这里的“真”由指定商映射的非平凡核来判断，不以 \(G/N\) 是否与 \(G\) 抽象同构来定义。

\ProseParagraphBreak
单群可以有许多真子群，定义排除的只是真正介于单位子群与整个群之间的正规子群。由\cref{b1:prop:normal-basic-examples}的核正规性，一个从单群出发的同态要么为平凡同态，要么为单射：若核为全群则像只有单位元，若核为单位子群，则\cref{b1:prop:kernel-fibers}保证单射。因此单群的每个同态像要么平凡，要么同构于原群；后一情形只需将同态的目标限制为实际像，便得到双射同态。

\begin{proposition}\label{b1:prop:abelian-simple}
设 \(G\) 为非平凡交换群。则 \(G\) 是单群，当且仅当 \(|G|\) 为素数。
\end{proposition}

\begin{proof}
若 \(\lvert G\rvert=p\) 为素数，则\cref{b1:thm:lagrange}使每个子群的阶只能为 \(1\) 或 \(p\)，所以没有非平凡真子群，尤其没有非平凡真正规子群。

反过来，设 \(G\) 是交换单群，取 \(g\ne1\)。
交换群的所有子群均正规，故非平凡子群 \(\langle g\rangle\) 必为 \(G\)。
如果 \(g\) 为无限阶，则 \(\langle g^2\rangle\) 非平凡且不含 \(g\)：否则 \(g=g^{2k}\) 给出非零整数次幂等于 \(1\)。这与单性矛盾。
故 \(g\) 有有限阶 \(n>1\)。若 \(n=ab\)，其中 \(a,b>1\)，则 \(g^a\ne1\)，而 \(\langle g^a\rangle\) 至多有 \(b<n\) 个元素，再次得到非平凡真子群。因此 \(n\) 为素数。
\end{proof}

非交换单群的构造需要比本章更多的工具；单群概念提出的结构问题是：能否通过非平凡真正规子群及相应商群研究原群？这种分解使用由原群乘法诱导的商群运算，因而必须以正规子群为依据。
